% ============================================================
% 纯答案册·测评卷 —— 工具/答案抽册器.py 抽册生成件（勿手改；第三档＝纯答案单独成册）
% 源件（只读）：工作区/M3-第2章量产0913/成卷/测评卷/main.tex（md5 fe9628bd94473d485829333fd64dff6b）
%% 口径：题面不重印；逐题＝题号(印面号同正册)＋[答案]＋[详解]，值/详解照源件逐字
%       （钉值门硬断言：册值≡正册件内值≡值台账值tex，逐字全等）。册序＝正册装配序＝印面号 1..19 升序（同序同号）；键序断言基准＝题面库冻结 manifest canonical 键序。
%% 三档导出：整体 main-true／纯题 main-false（在产）｜本册＝纯答案册（本件）
%% 双档：ansbook-true.tex＝详解印本；ansbook-false.tex＝纯值速查本（详解不印）
% 编译：xelatex <壳> ×2，cwd＝本目录；sty＝qp-m3.sty 本地挂载（md5 7c3930362be8a0a2bdf21bbf8ac16573，锁校验）
% ============================================================
\documentclass[fontset=none]{ctexart}
\usepackage{qp-m3}
% —— 印面逐字符号路由补钉（承源件；− 走 TNR、▱ NSC 子块；禁改 sty 保 md5） ——
\xeCJKDeclareCharClass{Default}{"2212}%
\setCJKfamilyfont{nscblk}[Path=C:/提示词/工作区/字替对照-0909/variantF/fonts/,BoldFont=NSC-w600.ttf]{NSC-w500.ttf}%
\xeCJKDeclareSubCJKBlock{nscblk}{"25B1}%
\newcommand{\pxparallelogram}{{\CJKfamily{nscblk}▱}}%
% —— 断行微调（承母版实测档，三零门承重；\emergencystretch 不另设） ——
\xeCJKsetup{CJKglue={\hskip 0.10em plus 0.08em minus 0.05em}}%
% —— 纯答案册全线括线（长详解可跨栏断，承练习件全线括线口径；灰底盒不可跨栏断） ——
\ansblockgrayfalse
% —— 双档开关（false 档＝纯值速查：答案印、详解不印；件面开关，不动 sty 本体） ——
\newif\ifansbookdetail
\ifdefined\ansbookdetail
  \ifnum\ansbookdetail=0 \ansbookdetailfalse\else\ansbookdetailtrue\fi
\else
  \ansbookdetailtrue
\fi
% —— 页脚件名（承源件章名；件名＝<件型>答案册） ——
\renewcommand{\qpzhangming}{}%
\renewcommand{\qpjianming}{答案册答案册}%

\begin{document}
\par\glueguard{1}\addvspace{2pt}\noindent\biaoqian{【纯答案册】}{\fangsong 题面不重印\quad 印面号与正册同号同序}\par\addvspace{2pt}

\begin{multicols}{2}
\emergencystretch=1em

\begin{ansblock}[2章-测-1]
% ans:2章-测-1
\ansitem{1}{A}
\ifansbookdetail
\fi
\end{ansblock}

\begin{ansblock}[2章-测-2]
% ans:2章-测-2
\ansitem{2}{C}
\ifansbookdetail
\fi
\end{ansblock}

\begin{ansblock}[2章-测-3]
% ans:2章-测-3
\ansitem{3}{A}
\ifansbookdetail
\fi
\end{ansblock}

\begin{ansblock}[2章-测-4]
% ans:2章-测-4
\ansitem{4}{C（两圆并取）}
\ifansbookdetail
\fi
\end{ansblock}

\begin{ansblock}[2章-测-5]
% ans:2章-测-5
\ansitem{5}{A}
\ifansbookdetail
\fi
\end{ansblock}

\begin{ansblock}[2章-测-6]
% ans:2章-测-6
\ansitem{6}{B（两条射线）}
\ifansbookdetail
\fi
\end{ansblock}

\begin{ansblock}[2章-测-7]
% ans:2章-测-7
\ansitem{7}{A（\(m=-\dfrac{1}{2}\)）}
\ifansbookdetail
\fi
\end{ansblock}

\begin{ansblock}[2章-测-8]
% ans:2章-测-8
\ansitem{8}{A（\([2\sqrt{3},4)\)）}
\ifansbookdetail
\fi
\end{ansblock}

\begin{ansblock}[2章-测-9]
% ans:2章-测-9
\ansitem{9}{AB}
\ifansbookdetail
\fi
\end{ansblock}

\begin{ansblock}[2章-测-10]
% ans:2章-测-10
\ansitem{10}{BCD}
\ifansbookdetail
\fi
\end{ansblock}

\begin{ansblock}[2章-测-11]
% ans:2章-测-11
\ansitem{11}{BCD}
\ifansbookdetail
\fi
\end{ansblock}

\begin{ansblock}[2章-测-12]
% ans:2章-测-12
\ansitem{12}{\(y=1\)}
\ifansbookdetail
\fi
\end{ansblock}

\begin{ansblock}[2章-测-13]
% ans:2章-测-13
\ansitem{13}{\((x-1)^{2}+y^{2}=4\)}
\ifansbookdetail
\fi
\end{ansblock}

\begin{ansblock}[2章-测-14]
% ans:2章-测-14
\ansitem{14}{\(1+\sqrt{3}\)}
\ifansbookdetail
\fi
\end{ansblock}

\begin{ansblock}[2章-测-15]
% ans:2章-测-15
\ansitem{15}{(1)\(y^{2}=4x\)；(2)直线 \(AB\) 恒过定点 \((-1,0)\)}
\ifansbookdetail
\fi
\end{ansblock}

\begin{ansblock}[2章-测-16]
% ans:2章-测-16
\ansitem{16}{(1)\(\dfrac{x^{2}}{6}+\dfrac{y^{2}}{3}=1\)；(2)\(x-\sqrt{7}y+\sqrt{3}=0\) 或 \(x+\sqrt{7}y+\sqrt{3}=0\)}
\ifansbookdetail
\fi
\end{ansblock}

\begin{ansblock}[2章-测-17]
% ans:2章-测-17
\ansitem{17}{\(-\dfrac{4}{3}\)}
\ifansbookdetail
\fi
\end{ansblock}

\begin{ansblock}[2章-测-18]
% ans:2章-测-18
\ansitem{18}{\(\dfrac{2\sqrt{3}}{3}\)}
\ifansbookdetail
\fi
\end{ansblock}

\begin{ansblock}[2章-测-19]
% ans:2章-测-19
\ansitem{19}{\(a=1\)}
\ifansbookdetail
\fi
\end{ansblock}

% ---- 尾块（\tailfill[30mm] 定高参——钉档 §三.3：无参在平衡栏呈「内容尾随框」，贴栏底须定高参；实测无参致末页平衡 Overfull\vbox 12.99pt，定高 30mm 三零） ----
\tailfill[30mm]

\end{multicols}
\end{document}
